\section{Four-Vectors and Covectors}
\label{sec:four_vectors_and_covectors}

A relativistic law should retain its form when one inertial observer is replaced by another. This requirement is easiest to express using objects with known Lorentz-transformation rules.

Four-vectors generalize ordinary spatial vectors to Minkowski spacetime. Covectors are the corresponding dual objects. The metric connects the two, and their contraction produces Lorentz scalars.

\subsection{Contravariant four-vectors}

A contravariant four-vector is an ordered set of four components

\[
A^{\mu}
=
(A^{0},A^{1},A^{2},A^{3})
\]

that transforms under a Lorentz transformation according to

\[
\boxed{
A'^{\mu}
=
\Lambda^{\mu}{}_{\nu}A^{\nu}
}.
\]

The superscript is not an exponent. It identifies the component and indicates the transformation type.

A collection of four numbers is not automatically a four-vector. The defining condition is the transformation law. Four quantities assembled into one column become a four-vector only if they transform together through the Lorentz matrix.

\subsection{Examples of contravariant four-vectors}

Important examples include:

\begin{itemize}
    \item the position four-vector
    \[
    x^{\mu}=(ct,\mathbf{x});
    \]

    \item the displacement four-vector
    \[
    \Delta x^{\mu}=x_{2}^{\mu}-x_{1}^{\mu};
    \]

    \item the four-velocity
    \[
    U^{\mu}=\frac{\mathrm{d}x^{\mu}}{\mathrm{d}\tau};
    \]

    \item the four-momentum
    \[
    p^{\mu}=mU^{\mu};
    \]

    \item a four-current
    \[
    J^{\mu};
    \]

    \item the electromagnetic four-potential
    \[
    A^{\mu}.
    \]
\end{itemize}

The last two examples will be developed in later chapters.

\subsection{Covectors}

A covector has components

\[
A_{\mu}
=
(A_{0},A_{1},A_{2},A_{3})
\]

and transforms with the inverse Lorentz transformation:

\[
\boxed{
A'_{\mu}
=
(\Lambda^{-1})^{\nu}{}_{\mu}A_{\nu}
}.
\]

Covectors are linear maps from vectors to scalars. Given a vector \(V^{\mu}\), the covector \(A_{\mu}\) produces the contraction

\[
A_{\mu}V^{\mu}.
\]

This quantity is invariant under Lorentz transformations.

In Euclidean Cartesian coordinates, vectors and covectors often have the same numerical components, so their distinction can remain hidden. In Minkowski spacetime, lowering a spatial index changes its sign, making the distinction immediately visible.

\subsection{Lowering an index}

The Minkowski metric converts a contravariant vector into a covector:

\[
\boxed{
A_{\mu}
=
\eta_{\mu\nu}A^{\nu}
}.
\]

For

\[
A^{\mu}
=
(A^{0},A^{1},A^{2},A^{3}),
\]

and

\[
\eta_{\mu\nu}
=
\operatorname{diag}(1,-1,-1,-1),
\]

we obtain

\[
\boxed{
A_{\mu}
=
(A^{0},-A^{1},-A^{2},-A^{3})
}.
\]

Thus

\[
A_{0}=A^{0},
\qquad
A_{i}=-A^{i}
\]

for spatial indices \(i=1,2,3\).

\subsection{Raising an index}

The inverse metric raises a covariant index:

\[
\boxed{
A^{\mu}
=
\eta^{\mu\nu}A_{\nu}
}.
\]

Since

\[
\eta^{\mu\nu}
=
\operatorname{diag}(1,-1,-1,-1),
\]

raising a spatial index changes its sign again:

\[
A^{0}=A_{0},
\qquad
A^{i}=-A_{i}.
\]

Applying lowering followed by raising returns the original vector:

\[
\eta^{\mu\alpha}\eta_{\alpha\nu}A^{\nu}
=
\delta^{\mu}{}_{\nu}A^{\nu}
=
A^{\mu}.
\]

\subsection{The Minkowski scalar product}

The scalar product of four-vectors \(A^{\mu}\) and \(B^{\mu}\) is

\[
\boxed{
A\cdot B
=
\eta_{\mu\nu}A^{\mu}B^{\nu}
=
A_{\mu}B^{\mu}
}.
\]

Expanding the components,

\[
\boxed{
A\cdot B
=
A^{0}B^{0}
-
A^{1}B^{1}
-
A^{2}B^{2}
-
A^{3}B^{3}
}.
\]

Using spatial vector notation,

\[
A\cdot B
=
A^{0}B^{0}
-
\mathbf{A}\cdot\mathbf{B}.
\]

The contraction is a Lorentz scalar. Its numerical value is the same in every inertial frame.

\subsection{The norm of a four-vector}

The squared norm is

\[
\boxed{
A^{2}
:=
A\cdot A
=
A_{\mu}A^{\mu}
}.
\]

Explicitly,

\[
A^{2}
=
(A^{0})^{2}
-
|\mathbf{A}|^{2}.
\]

Unlike a Euclidean norm, this quantity can be positive, zero, or negative.

A four-vector is called

\begin{itemize}
    \item timelike if \(A^{2}>0\);
    \item null if \(A^{2}=0\);
    \item spacelike if \(A^{2}<0\).
\end{itemize}

For the displacement four-vector, this classification reproduces the causal classification of spacetime separations.

\subsection{Worked example: lowering a vector index}

Let

\[
A^{\mu}=(5,2,-1,3).
\]

Then

\[
A_{\mu}
=
\eta_{\mu\nu}A^{\nu}
=
(5,-2,1,-3).
\]

The squared norm is

\[
\begin{aligned}
A_{\mu}A^{\mu}
&=
5(5)
+(-2)(2)
+(1)(-1)
+(-3)(3)\\
&=
25-4-1-9\\
&=
11.
\end{aligned}
\]

Equivalently,

\[
A^{2}
=
5^{2}-2^{2}-(-1)^{2}-3^{2}
=11.
\]

Since

\[
A^{2}>0,
\]

the vector is timelike.

\subsection{Worked example: scalar product}

Let

\[
A^{\mu}=(4,1,2,0),
\qquad
B^{\mu}=(3,-2,1,5).
\]

Then

\[
\begin{aligned}
A\cdot B
&=
A^{0}B^{0}
-A^{1}B^{1}
-A^{2}B^{2}
-A^{3}B^{3}\\
&=
4(3)-1(-2)-2(1)-0(5)\\
&=
12+2-2\\
&=
12.
\end{aligned}
\]

Using the lowered vector

\[
A_{\mu}=(4,-1,-2,0),
\]

we obtain the same result:

\[
A_{\mu}B^{\mu}
=
4(3)+(-1)(-2)+(-2)(1)+0(5)
=12.
\]

\subsection{The position four-vector}

The position four-vector is

\[
\boxed{
x^{\mu}=(ct,\mathbf{x})
}.
\]

Its covariant components are

\[
\boxed{
x_{\mu}=(ct,-\mathbf{x})
}.
\]

The contraction

\[
x_{\mu}x^{\mu}
=
c^{2}t^{2}-|\mathbf{x}|^{2}
\]

is the squared interval from the coordinate origin to the event.

The coordinate origin itself has no physical privilege. Translating the origin changes \(x^{\mu}\), but differences

\[
\Delta x^{\mu}
\]

between events transform as four-vectors and carry direct geometric meaning.

\subsection{The four-velocity}

For a massive particle, the four-velocity is defined using proper time:

\[
\boxed{
U^{\mu}
=
\frac{\mathrm{d}x^{\mu}}{\mathrm{d}\tau}
}.
\]

Since

\[
\frac{\mathrm{d}t}{\mathrm{d}\tau}
=
\gamma,
\]

we obtain

\[
\begin{aligned}
U^{0}
&=
\frac{\mathrm{d}(ct)}{\mathrm{d}\tau}
=
\gamma c,\\
\mathbf{U}
&=
\frac{\mathrm{d}\mathbf{x}}{\mathrm{d}\tau}
=
\gamma\frac{\mathrm{d}\mathbf{x}}{\mathrm{d}t}
=
\gamma\mathbf{v}.
\end{aligned}
\]

Therefore

\[
\boxed{
U^{\mu}
=
\gamma(c,\mathbf{v})
}.
\]

The four-velocity is not simply \((c,\mathbf{v})\). The factor \(\gamma\) is required for the object to transform as a four-vector.

\subsection{Normalization of the four-velocity}

The squared norm is

\[
\begin{aligned}
U_{\mu}U^{\mu}
&=
\gamma^{2}
(c^{2}-v^{2})\\
&=
\gamma^{2}c^{2}
\left(1-
\frac{v^{2}}{c^{2}}
\right)\\
&=
c^{2}.
\end{aligned}
\]

Thus

\[
\boxed{
U_{\mu}U^{\mu}=c^{2}
}.
\]

Every massive particle has the same four-velocity norm, even though different observers assign different three-velocities to it.

This provides an important structural check. Any proposed formula for a four-velocity should satisfy this normalization.

\subsection{The four-momentum}

For a particle of invariant mass \(m\), define

\[
\boxed{
p^{\mu}=mU^{\mu}
}.
\]

Using the four-velocity components,

\[
p^{\mu}
=
(\gamma mc,\gamma m\mathbf{v}).
\]

The relativistic three-momentum is

\[
\mathbf{p}
=
\gamma m\mathbf{v},
\]

and the relativistic energy is

\[
E
=
\gamma mc^{2}.
\]

Therefore

\[
\boxed{
p^{\mu}
=
\left(
\frac{E}{c},
\mathbf{p}
\right)
}.
\]

Lowering the index gives

\[
\boxed{
p_{\mu}
=
\left(
\frac{E}{c},
-\mathbf{p}
\right)
}.
\]

\subsection{Norm of the four-momentum}

Since

\[
p^{\mu}=mU^{\mu},
\]

we have

\[
\begin{aligned}
p_{\mu}p^{\mu}
&=
m^{2}U_{\mu}U^{\mu}\\
&=
m^{2}c^{2}.
\end{aligned}
\]

Thus

\[
\boxed{
p_{\mu}p^{\mu}=m^{2}c^{2}
}.
\]

Using

\[
p^{\mu}=\left(E/c,\mathbf{p}\right),
\]

this becomes

\[
\frac{E^{2}}{c^{2}}
-|
\mathbf{p}|^{2}
=
m^{2}c^{2}.
\]

Multiplying by \(c^{2}\),

\[
\boxed{
E^{2}
=
|
\mathbf{p}|^{2}c^{2}
+
m^{2}c^{4}
}.
\]

This relation will be revisited in the section on relativistic invariants.

\subsection{Massless four-momentum}

For a massless particle,

\[
m=0,
\]

so

\[
p_{\mu}p^{\mu}=0.
\]

The energy-momentum relation reduces to

\[
E^{2}=|
\mathbf{p}|^{2}c^{2},
\]

or, for positive energy,

\[
\boxed{
E=c|
\mathbf{p}|
}.
\]

A massless particle has null four-momentum. Its four-momentum is well defined even though proper time and four-velocity are not defined along a null worldline in the same way as for a massive particle.

\subsection{Basis-vector interpretation}

A four-vector is a geometric object independent of coordinates. In a basis \(\{\mathbf{e}_{\mu}\}\), it may be written as

\[
\mathbf{A}
=
A^{\mu}\mathbf{e}_{\mu}.
\]

When the inertial frame changes, both the components and the basis change in compensating ways, leaving the geometric vector unchanged.

A covector belongs to the dual basis \(\{\boldsymbol{\theta}^{\mu}\}\):

\[
\boldsymbol{\alpha}
=
\alpha_{\mu}\boldsymbol{\theta}^{\mu}.
\]

The dual pairing satisfies

\[
\boldsymbol{\theta}^{\mu}(\mathbf{e}_{\nu})
=
\delta^{\mu}{}_{\nu}.
\]

This geometric viewpoint explains why upper and lower indices transform differently.

\subsection{Structural checks}

A proposed four-vector calculation should satisfy:

\begin{enumerate}
    \item the object has four components;
    \item its components transform with the Lorentz matrix;
    \item lowering or raising an index uses the metric exactly once;
    \item scalar products contain one upper and one lower occurrence of the contracted index;
    \item the norm has the correct sign for its causal type;
    \item the four-velocity satisfies \(U^{2}=c^{2}\);
    \item the four-momentum satisfies \(p^{2}=m^{2}c^{2}\).
\end{enumerate}

\subsection{Common mistakes}

Typical errors include:

\begin{enumerate}
    \item treating the index as an exponent;
    \item assuming that every list of four numbers is a four-vector;
    \item lowering an index without changing the spatial signs;
    \item using the Euclidean scalar product instead of the Minkowski scalar product;
    \item writing \(U^{\mu}=(c,\mathbf{v})\) and omitting \(\gamma\);
    \item confusing the invariant mass \(m\) with the frame-dependent energy \(E\);
    \item defining a four-velocity for a light ray using proper time, even though \(\mathrm{d}\tau=0\) on a null worldline.
\end{enumerate}

\subsection{What has been established}

A contravariant four-vector transforms as

\[
A'^{\mu}
=
\Lambda^{\mu}{}_{\nu}A^{\nu},
\]

while the metric produces the corresponding covector

\[
A_{\mu}
=
\eta_{\mu\nu}A^{\nu}.
\]

Their contraction defines the Lorentz-invariant scalar product. The position, four-velocity, and four-momentum are central examples:

\[
x^{\mu}=(ct,\mathbf{x}),
\]

\[
U^{\mu}=\gamma(c,\mathbf{v}),
\]

and

\[
p^{\mu}=(E/c,\mathbf{p}).
\]

The next section develops the index notation needed to manipulate these objects reliably, including free indices, dummy indices, contractions, and tensor equations.